% GENERATED FILE. Edit canonical lessons, then run npm run generate:books.
% canonical-source-hash: fnv1a64:b4930aa33ca2a8a8
% canonical-lessons: GU-C371-phaydo, GU-C371-labh, GU-C371-upyog, GU-C371-karaj, GU-C371-khovu

\chapter{Gains and Losses}
\label{ch:gu-gains-and-losses}
\hlchaptermodality{371}

\begin{chapteropening}
\textbf{By the end of this chapter:} \emph{I can weigh what we gain and what we lose.}

Complete the last lesson of chapter 371: I can weigh what we gain and what we lose.
\end{chapteropening}

\section[\texorpdfstring{\gu{ફાયદો} (phāydo)}{ફાયદો (phāydo)}]{\gu{ફાયદો} (phāydo) — a benefit}

\label{lesson:GU-C371-phaydo}



\begin{quote}
Before the new one: say the Gujarati for to spoil, then the Gujarati for to put up with.
\end{quote}

\subsection*{You'll want to know: \gu{ફાયદો}}
\textbf{\gu{ફાયદો}} — \emph{phāydo} — \textquotedblleft{}a benefit\textquotedblright{}.

\begin{quote}
\gu{જો} \gu{તમે} \gu{વહેલા} \gu{ટિકિટ} \gu{ખરીદશો} \gu{તો} \gu{ફાયદો} \gu{થશે}.
\end{quote}

\begin{quote}
\emph{jo tame vahelā ṭikiṭ kharīdśo to phāydo thaśe.}
\end{quote}

\begin{quote}
If you buy the tickets early, you will save.
\end{quote}

A masculine noun.

\begin{cousinweb}
\textbf{\gu{ફાયદો}} (\emph{phāydo}) is usually traced to Arabic \emph{fāʾida}, \textquotedblleft{}a benefit, a profit\textquotedblright{}, through Persian.
\end{cousinweb}

\begin{grammarlens}[title={What you've built}]
One more word for saying what happens if.
\end{grammarlens}

\subsection*{Your turn}
\begin{itemize}
\raggedright
  \item \emph{Say it:} \emph{phāydo}
  \item \emph{Say it:} \emph{phāydo}, once more
  \item \emph{Say it:} \emph{jo tame vahelā ṭikiṭ kharīdśo to phāydo thaśe.} — If you buy the tickets early, you will save.
  \item \emph{Recall:} say the Gujarati for to spoil, then the Gujarati for to put up with, then say \emph{phāydo} again
\end{itemize}

\begin{tcolorbox}[breakable,colback=teal!4,colframe=teal!35!black,title={Before you move on}]
What does \gu{ફાયદો} mean? (\textquotedblleft{}A benefit\textquotedblright{}.) Say it once more.
\end{tcolorbox}
\section[\texorpdfstring{\gu{લાભ} (lābh)}{લાભ (lābh)}]{\gu{લાભ} (lābh) — an advantage}

\label{lesson:GU-C371-labh}



\begin{quote}
Before the new one: say the Gujarati for to put up with, then the Gujarati for a benefit.
\end{quote}

\subsection*{You'll want to know: \gu{લાભ}}
\textbf{\gu{લાભ}} — \emph{lābh} — \textquotedblleft{}an advantage\textquotedblright{}.

\begin{quote}
\gu{જો} \gu{તમે} \gu{આ} \gu{યોજનામાં} \gu{જોડાશો} \gu{તો} \gu{તમને} \gu{ઘણો} \gu{લાભ} \gu{થશે}.
\end{quote}

\begin{quote}
\emph{jo tame ā yojnāmā̃ joḍāśo to tamne ghaṇo lābh thaśe.}
\end{quote}

\begin{quote}
If you join this scheme, you will gain a lot.
\end{quote}

A masculine noun.

\begin{cousinweb}
\textbf{\gu{લાભ}} (\emph{lābh}) is a tatsama, borrowed from Sanskrit \emph{lābha}, \textquotedblleft{}finding, gain\textquotedblright{} (Monier-Williams).
\end{cousinweb}

\begin{grammarlens}[title={What you've built}]
One more word for saying what happens if.
\end{grammarlens}

\subsection*{Your turn}
\begin{itemize}
\raggedright
  \item \emph{Say it:} \emph{lābh}
  \item \emph{Say it:} \emph{lābh}, once more
  \item \emph{Say it:} \emph{jo tame ā yojnāmā̃ joḍāśo to tamne ghaṇo lābh thaśe.} — If you join this scheme, you will gain a lot.
  \item \emph{Recall:} say the Gujarati for to put up with, then the Gujarati for a benefit, then say \emph{lābh} again
\end{itemize}

\begin{tcolorbox}[breakable,colback=teal!4,colframe=teal!35!black,title={Before you move on}]
What does \gu{લાભ} mean? (\textquotedblleft{}An advantage\textquotedblright{}.) Say it once more.
\end{tcolorbox}
\section[\texorpdfstring{\gu{ઉપયોગ} (upyog)}{ઉપયોગ (upyog)}]{\gu{ઉપયોગ} (upyog) — use; usefulness}

\label{lesson:GU-C371-upyog}



\begin{quote}
Before the new one: say the Gujarati for a benefit, then the Gujarati for an advantage.
\end{quote}

\subsection*{You'll want to know: \gu{ઉપયોગ}}
\textbf{\gu{ઉપયોગ}} — \emph{upyog} — \textquotedblleft{}use; usefulness\textquotedblright{}.

\begin{quote}
\gu{જો} \gu{તમે} \gu{આ} \gu{પુસ્તકનો} \gu{ઉપયોગ} \gu{કરશો} \gu{તો} \gu{પાઠ} \gu{સહેલો} \gu{થશે}.
\end{quote}

\begin{quote}
\emph{jo tame ā pustakno upyog karśo to pāṭh sahelo thaśe.}
\end{quote}

\begin{quote}
If you use this book, the lesson will be easy.
\end{quote}

A masculine noun.

\begin{cousinweb}
\textbf{\gu{ઉપયોગ}} (\emph{upyog}) is a tatsama, borrowed from Sanskrit \emph{upayoga}, \textquotedblleft{}employment, use, application\textquotedblright{} (Monier-Williams).
\end{cousinweb}

\begin{grammarlens}[title={What you've built}]
One more word for saying what happens if.
\end{grammarlens}

\subsection*{Your turn}
\begin{itemize}
\raggedright
  \item \emph{Say it:} \emph{upyog}
  \item \emph{Say it:} \emph{upyog}, once more
  \item \emph{Say it:} \emph{jo tame ā pustakno upyog karśo to pāṭh sahelo thaśe.} — If you use this book, the lesson will be easy.
  \item \emph{Recall:} say the Gujarati for a benefit, then the Gujarati for an advantage, then say \emph{upyog} again
\end{itemize}

\begin{tcolorbox}[breakable,colback=teal!4,colframe=teal!35!black,title={Before you move on}]
What does \gu{ઉપયોગ} mean? (\textquotedblleft{}Use; usefulness\textquotedblright{}.) Say it once more.
\end{tcolorbox}
\section[\texorpdfstring{\gu{કરજ} (karaj)}{કરજ (karaj)}]{\gu{કરજ} (karaj) — a debt}

\label{lesson:GU-C371-karaj}



\begin{quote}
Before the new one: say the Gujarati for an advantage, then the Gujarati for use.
\end{quote}

\subsection*{You'll want to know: \gu{કરજ}}
\textbf{\gu{કરજ}} — \emph{karaj} — \textquotedblleft{}a debt\textquotedblright{}.

\begin{quote}
\gu{જો} \gu{તમે} \gu{વધારે} \gu{કરજ} \gu{લેશો} \gu{તો} \gu{વ્યાજ} \gu{પણ} \gu{વધારે} \gu{ભરવું} \gu{પડશે}.
\end{quote}

\begin{quote}
\emph{jo tame vadhāre karaj leśo to vyāj paṇ vadhāre bharvũ paḍśe.}
\end{quote}

\begin{quote}
If you take on more debt, you will have to pay more interest too.
\end{quote}

A neuter noun.

\begin{cousinweb}
\textbf{\gu{કરજ}} (\emph{karaj}) is usually traced to Arabic \emph{qarḍ}, \textquotedblleft{}a loan\textquotedblright{}, through Persian \emph{qarz}.
\end{cousinweb}

\begin{grammarlens}[title={What you've built}]
One more word for saying what happens if.
\end{grammarlens}

\subsection*{Your turn}
\begin{itemize}
\raggedright
  \item \emph{Say it:} \emph{karaj}
  \item \emph{Say it:} \emph{karaj}, once more
  \item \emph{Say it:} \emph{jo tame vadhāre karaj leśo to vyāj paṇ vadhāre bharvũ paḍśe.} — If you take on more debt, you will have to pay more interest too.
  \item \emph{Recall:} say the Gujarati for an advantage, then the Gujarati for use, then say \emph{karaj} again
\end{itemize}

\begin{tcolorbox}[breakable,colback=teal!4,colframe=teal!35!black,title={Before you move on}]
What does \gu{કરજ} mean? (\textquotedblleft{}A debt\textquotedblright{}.) Say it once more.
\end{tcolorbox}
\section[\texorpdfstring{\gu{ખોવું} (khovũ)}{ખોવું (khovũ)}]{\gu{ખોવું} (khovũ) — to lose for good}

\label{lesson:GU-C371-khovu}



\begin{quote}
Before the new one: say the Gujarati for use, then the Gujarati for a debt.
\end{quote}

\subsection*{You'll want to know: \gu{ખોવું}}
\textbf{\gu{ખોવું}} — \emph{khovũ} — \textquotedblleft{}to lose for good\textquotedblright{}.

\begin{quote}
\gu{જો} \gu{તમે} \gu{આ} \gu{તક} \gu{ખોશો} \gu{તો} \gu{ફરી} \gu{તક} \gu{નહીં} \gu{આવે}.
\end{quote}

\begin{quote}
\emph{jo tame ā tak khośo to pharī tak nahī̃ āve.}
\end{quote}

\begin{quote}
If you lose this chance, it will not come again.
\end{quote}

Past, with \gu{તેણે} (he or she …): \gu{તેણે} \gu{પોતાનું} \gu{પાકીટ} \gu{ખોયું}, \emph{teṇe potānũ pākīṭ khoyũ}.

\begin{cousinweb}
\textbf{\gu{ખોવું}} (\emph{khovũ}) is inherited from Sanskrit \emph{kṣapayati}, \textquotedblleft{}destroys\textquotedblright{}; CDIAL lists the Gujarati verb, \textquotedblleft{}to lose, to waste\textquotedblright{} (CDIAL 3651). \textbf{\gu{ખોવાવું}} (\emph{khovāvũ}), \textquotedblleft{}to get lost\textquotedblright{}, is its partner.
\end{cousinweb}

\begin{grammarlens}[title={What you've built}]
One more word for saying what happens if.
\end{grammarlens}

\subsection*{Your turn}
\begin{itemize}
\raggedright
  \item \emph{Say it:} \emph{khovũ}
  \item \emph{Say it:} \emph{khovũ}, once more
  \item \emph{Say it:} \emph{jo tame ā tak khośo to pharī tak nahī̃ āve.} — If you lose this chance, it will not come again.
  \item \emph{Recall:} say the Gujarati for use, then the Gujarati for a debt, then say \emph{khovũ} again
\end{itemize}

\begin{tcolorbox}[breakable,colback=teal!4,colframe=teal!35!black,title={Before you move on}]
What does \gu{ખોવું} mean? (\textquotedblleft{}To lose for good\textquotedblright{}.) Say it once more.
\end{tcolorbox}
